\submittedby{Manas Singh}

Let $X=\{a,b,c\}$. The group $S_3$ acts on $\T_X$ by $f\cdot\T=\{f(U)\mid U\in\T\}$. Every
topology contains $\emptyset$ and $X$, so we list all $29$ topologies explicitly, grouped by
orbit. Within each orbit, the topologies are obtained from one another by permuting $a,b,c$.

\begin{enumerate}
  \item \textbf{Indiscrete topology} (orbit size $1$):
  \begin{itemize}
    \item $\{\emptyset,X\}$
  \end{itemize}

  \item \textbf{One non-trivial open set, a singleton} (orbit size $3$):
  \begin{itemize}
    \item $\{\emptyset,\{a\},X\}$
    \item $\{\emptyset,\{b\},X\}$
    \item $\{\emptyset,\{c\},X\}$
  \end{itemize}

  \item \textbf{One non-trivial open set, a two-element set} (orbit size $3$):
  \begin{itemize}
    \item $\{\emptyset,\{a,b\},X\}$
    \item $\{\emptyset,\{b,c\},X\}$
    \item $\{\emptyset,\{a,c\},X\}$
  \end{itemize}

  \item \textbf{Partition topology} (orbit size $3$):
  \begin{itemize}
    \item $\{\emptyset,\{a\},\{b,c\},X\}$
    \item $\{\emptyset,\{b\},\{a,c\},X\}$
    \item $\{\emptyset,\{c\},\{a,b\},X\}$
  \end{itemize}

  \item \textbf{Nested chain $\{x\}\subset\{x,y\}\subset X$} (orbit size $6$, one for each
  ordered pair $(x,y)$ of distinct points):
  \begin{itemize}
    \item $\{\emptyset,\{a\},\{a,b\},X\}$
    \item $\{\emptyset,\{a\},\{a,c\},X\}$
    \item $\{\emptyset,\{b\},\{a,b\},X\}$
    \item $\{\emptyset,\{b\},\{b,c\},X\}$
    \item $\{\emptyset,\{c\},\{a,c\},X\}$
    \item $\{\emptyset,\{c\},\{b,c\},X\}$
  \end{itemize}

  \item \textbf{Two disjoint singletons and their union} (orbit size $3$, determined by the
  excluded point):
  \begin{itemize}
    \item $\{\emptyset,\{a\},\{b\},\{a,b\},X\}$
    \item $\{\emptyset,\{b\},\{c\},\{b,c\},X\}$
    \item $\{\emptyset,\{a\},\{c\},\{a,c\},X\}$
  \end{itemize}

  \item \textbf{One singleton with both two-element sets containing it} (orbit size $3$,
  determined by the singleton):
  \begin{itemize}
    \item $\{\emptyset,\{a\},\{a,b\},\{a,c\},X\}$
    \item $\{\emptyset,\{b\},\{a,b\},\{b,c\},X\}$
    \item $\{\emptyset,\{c\},\{a,c\},\{b,c\},X\}$
  \end{itemize}

  \item \textbf{Two singletons $\{x\},\{y\}$ with $\{x,y\}$ and $\{x,z\}$} (orbit size $6$, one
  for each ordered pair $(x,y)$ of distinct points, where $z$ is the remaining point):
  \begin{itemize}
    \item $\{\emptyset,\{a\},\{b\},\{a,b\},\{a,c\},X\}$
    \item $\{\emptyset,\{b\},\{a\},\{a,b\},\{b,c\},X\}$
    \item $\{\emptyset,\{a\},\{c\},\{a,c\},\{a,b\},X\}$
    \item $\{\emptyset,\{c\},\{a\},\{a,c\},\{b,c\},X\}$
    \item $\{\emptyset,\{b\},\{c\},\{b,c\},\{a,b\},X\}$
    \item $\{\emptyset,\{c\},\{b\},\{b,c\},\{a,c\},X\}$
  \end{itemize}

  \item \textbf{Discrete topology} (orbit size $1$):
  \begin{itemize}
    \item $\Pow(X)=\{\emptyset,\{a\},\{b\},\{c\},\{a,b\},\{b,c\},\{a,c\},X\}$
  \end{itemize}
\end{enumerate}

Each listed collection contains $\emptyset$ and $X$ and is closed under unions and
intersections, so each is a topology. Types 1--9 are pairwise non-homeomorphic: they differ in
the number of open sets ($2,3,3,4,4,5,5,6,8$) or, when these agree, in their structure (types 2
and 3 differ in the size of the non-trivial open set; types 4 and 5 differ because the open sets
of type 4 are disjoint while those of type 5 are nested; types 6 and 7 differ because type 6 has
two singleton open sets and type 7 has only one). Within each type, permuting the points moves
any topology to any other, so each type is exactly one orbit.

Summing the orbit sizes gives
\[ 1+3+3+3+6+3+3+6+1=29, \]
which accounts for all $29$ topologies. Hence the number of orbits of $S_3$ on $\T_X$, i.e.\ the
number of homeomorphism classes of topologies on a three-point set, is exactly \textbf{9}.

\submittedby{Sarthak}

\medskip\noindent\textbf{All $29$ topologies.}\par\medskip

\textbf{Orbit 1: Indiscrete topology}
\[ \{\emptyset,X\}. \]

\textbf{Orbit 2: One singleton open}
\[ \{\emptyset,\{a\},X\},\qquad\{\emptyset,\{b\},X\},\qquad\{\emptyset,\{c\},X\}. \]

\textbf{Orbit 3: One doubleton open}
\[ \{\emptyset,\{a,b\},X\},\qquad\{\emptyset,\{a,c\},X\},\qquad\{\emptyset,\{b,c\},X\}. \]

\textbf{Orbit 4: One singleton and a doubleton containing it}
\begin{gather*}
\{\emptyset,\{a\},\{a,b\},X\},\qquad\{\emptyset,\{a\},\{a,c\},X\},\\
\{\emptyset,\{b\},\{a,b\},X\},\qquad\{\emptyset,\{b\},\{b,c\},X\},\\
\{\emptyset,\{c\},\{a,c\},X\},\qquad\{\emptyset,\{c\},\{b,c\},X\}.
\end{gather*}

\textbf{Orbit 5: One singleton and the complementary doubleton}
\[ \{\emptyset,\{a\},\{b,c\},X\},\qquad\{\emptyset,\{b\},\{a,c\},X\},\qquad
   \{\emptyset,\{c\},\{a,b\},X\}. \]

\textbf{Orbit 6: Two singletons and their union}
\begin{gather*}
\{\emptyset,\{a\},\{b\},\{a,b\},X\},\\
\{\emptyset,\{a\},\{c\},\{a,c\},X\},\\
\{\emptyset,\{b\},\{c\},\{b,c\},X\}.
\end{gather*}

\textbf{Orbit 7: One singleton and two doubletons containing it}
\begin{gather*}
\{\emptyset,\{a\},\{a,b\},\{a,c\},X\},\\
\{\emptyset,\{b\},\{a,b\},\{b,c\},X\},\\
\{\emptyset,\{c\},\{a,c\},\{b,c\},X\}.
\end{gather*}

\textbf{Orbit 8: Two singletons and three doubletons, omitting one doubleton}
\begin{gather*}
\{\emptyset,\{a\},\{b\},\{a,b\},\{a,c\},X\},\\
\{\emptyset,\{a\},\{c\},\{a,b\},\{a,c\},X\},\\
\{\emptyset,\{a\},\{b\},\{a,b\},\{b,c\},X\},\\
\{\emptyset,\{b\},\{c\},\{a,b\},\{b,c\},X\},\\
\{\emptyset,\{a\},\{c\},\{a,c\},\{b,c\},X\},\\
\{\emptyset,\{b\},\{c\},\{a,c\},\{b,c\},X\}.
\end{gather*}

\textbf{Orbit 9: Discrete topology}
\[ \{\emptyset,\{a\},\{b\},\{c\},\{a,b\},\{a,c\},\{b,c\},X\}. \]

Thus, the total number of topologies is
\[ 1+3+3+6+3+3+3+6+1=29. \]
Therefore, the number of orbits of $S_3$ acting on $\T_X$ is
\[ |\T_X/S_3|=9. \]

\paragraph{Explanation of orbit sizes.} The indiscrete topology and the discrete topology are
fixed by every permutation of $a,b,c$. Hence each forms an orbit of size $1$.

% REVIEW: never argues that topologies in different groups lie in different orbits (i.e. are not homeomorphic), which the count of 9 needs; the orbit-size remark below is only a heuristic.
The remaining seven orbits have sizes $3$ or $6$, depending on whether the representative
topology has nontrivial symmetry under permutations of $a,b,c$. Since
\[ 1+3+3+6+3+3+3+6+1=29, \]
the orbit decomposition accounts for all $29$ topologies on $X$.
